%% ==========================================================================
%%  Approach 10 --- minimum-spread families.  OPEN, and the best lead.
%%
%%  Source: RESIDUAL.md (2026-08-10) and update_48/.
%%
%%  NOTE ON PLACEMENT.  This is deliberately NOT part of \Cref{sec:approach9}.
%%  Approach 9 is the conditioned-remainder induction, whose state is a partial
%%  allocation under contracted costs; the technique below selects a PARTITION
%%  by a global optimisation and assigns it by a matching.  The two share
%%  nothing but the target, so they are kept apart.
%%
%%  STATUS.  The rule (F5*) is unrefuted over 2,574 instances and implies
%%  Conjecture 2.  At n = 3 everything here is proved except Conjecture G,
%%  which alone would close n = 3.
%% ==========================================================================
\section{Approach 10: Minimum-Spread Families}
\label{sec:approach10}

Every approach so far has selected an \emph{allocation}. This one selects a
\emph{family} --- an unassigned partition --- by a global optimisation, and only
then hands the bundles out. The optimisation is the obvious one suggested by
\Cref{sec:balancedclass}: Theorem~\ref{thm:balanced-class} succeeds exactly when
some family has every agent valuing all bundles within one unit, so the natural
move when no such family exists is to ask for the family that comes closest.

\begin{definition}[Spread]
\label{def:spread}
For a family $B = (B_1,\dots,B_n)$ --- a partition of $\items$ into $n$ possibly
empty bundles --- and an agent $i$, put
\[
  \mathrm{sp}_i(B) \;:=\; \max_t \cost_i(B_t) - \min_t \cost_i(B_t),
  \qquad
  \Sigma(B) \;:=\; \sum_{i \in \agents} \mathrm{sp}_i(B).
\]
A family is uniformly balanced (\Cref{def:unifbal}) exactly when
$\mathrm{sp}_i(B) \le 1$ for every $i$.
\end{definition}

\begin{conjecture}[The minimum-spread rule --- \textbf{open}]
\label{conj:f5}
Let $B$ minimise $\Sigma$ over all families and let $\sigma$ minimise
$\sum_i \cost_i(B_{\sigma(i)})$. Then $\alloc$ with $\bundle{i} =
B_{\sigma(i)}$ satisfies $\pathw{\alloc} \le 1$.
\end{conjecture}

Conjecture~\ref{conj:f5} implies Conjecture~\ref{conj:main} at once, and a
minimiser exists by definition, so unlike Conjecture~\ref{conj:balance-rule} it
always applies and needs no schedule. It is unrefuted over $2{,}574$ instances at
$n \le 6$, $m \le 8$ across eleven generators, in the strong form --- \emph{every}
minimiser works, not merely some --- on a corpus not selected for it
(\texttt{update\_48/minsum\_stress.py}).

\begin{remark}[What a proof may not use]
\label{rem:f5-closures}
Three mechanisms are already closed and none of the arguments below uses any of
them. First, the arc bound: Proposition~\ref{prop:arc-bound-blind} shows that on
the instance of Proposition~\ref{prop:no-balance} \emph{every} good allocation
has an arc of weight $-2$ or less, so no criterion of the form
``all arcs $\ge -1$'' --- that is, Corollary~\ref{cor:onestep} --- can be the
mechanism, and Remark~\ref{rem:arc-vs-path} draws the conclusion that any correct
certificate is path-based. Second, balance: over $42{,}711$ families of spread at
most $2$, requiring bundle sizes to differ by at most one admits $792$ families
whose matching is bad, so the balance signal of Remark~\ref{rem:balance} is not
sufficient here. Third, single-item exchange: $\cost(S) = \min(\max(0,\abs{S}-1),1)$
on a three-element set has $\cost = 1$ with every single removal leaving it at
$1$, so no step may move one item and expect a cost to move.
\end{remark}

% --------------------------------------------------------------------------
\subsection{Normalisation}
\label{sec:f5-normal}

Fix a family $B$ and put
$v_i(t) := \cost_i(B_t) - \min_s \cost_i(B_s) \ge 0$, so each $v_i$ attains the
value $0$. Three identities make the whole question finite for a fixed family:
\[
  \mathrm{sp}_i(B) = \max_t v_i(t), \qquad
  \edgew{\alloc}(i,k) = v_i(\sigma(i)) - v_i(\sigma(k)), \qquad
  \sum_i \cost_i(B_{\sigma(i)}) = \sum_i v_i(\sigma(i)) + \text{const}.
\]
The last one says a minimum-cost assignment is exactly one minimising
$F(\sigma) := \sum_i v_i(\sigma(i))$. Write $x_i := v_i(\sigma(i))$ for what
agent $i$ pays after normalisation.

\begin{observation}
\label{obs:f5-arcs}
The heaviest arc out of $i$ equals $x_i$ if $x_i \ge 1$, and is at most $0$ if
$x_i = 0$.
\end{observation}

\begin{proof}
$\max_{k \ne i}\edgew{\alloc}(i,k) = x_i - \min_{k \ne i} v_i(\sigma(k))$. If
$x_i \ge 1$ then $v_i$ does not vanish at $\sigma(i)$, so its zero lies at some
$\sigma(k)$ with $k \ne i$ and the inner minimum is $0$. If $x_i = 0$ then every
arc out of $i$ is $-v_i(\sigma(k)) \le 0$.
\end{proof}

A minimum-cost assignment maximises welfare over reassignments, so
$\alloc$ is envy-freeable by Theorem~\ref{thm:hs-characterisation} and
Proposition~\ref{prop:n3} applies at $n = 3$. Combining it with
Observation~\ref{obs:f5-arcs}:

\begin{proposition}[Criterion at $n = 3$]
\label{prop:f5-n3crit}
Let $n = 3$. Then $\alloc$ is good if and only if
\begin{enumerate}
  \item[(a)] $x_i \le 1$ for every $i$, and
  \item[(b)] there are no distinct $i,j,k$ with $x_i = x_j = 1$,
    $v_i(\sigma(j)) = 0$ and $v_j(\sigma(k)) = 0$.
\end{enumerate}
\end{proposition}

\begin{proof}
By Observation~\ref{obs:f5-arcs}, (a) is exactly ``every arc has weight at most
$1$''. Given (a), a two-path $i \to j \to k$ has weight $2$ if and only if both
its arcs weigh $1$, that is $x_i - v_i(\sigma(j)) = 1$ and
$x_j - v_j(\sigma(k)) = 1$; under (a) this is (b). Proposition~\ref{prop:n3}
then gives the claim.
\end{proof}

% --------------------------------------------------------------------------
\subsection{The main lemma}
\label{sec:f5-main}

\begin{theorem}[Small total spread is enough]
\label{thm:f5-sigma3}
Let $n = 3$ and let $\sigma$ be any minimum-cost assignment of a family $B$.
\begin{enumerate}
  \item[(i)] If some arc of $\envygraph{\alloc}$ has weight at least $2$, then
    $\Sigma(B) \ge 4$.
  \item[(ii)] If every arc has weight at most $1$ but some two-path has weight
    $2$, then $\Sigma(B) \ge 5$.
\end{enumerate}
In particular, if $\Sigma(B) \le 3$ then \emph{every} minimum-cost assignment of
$B$ is good.
\end{theorem}

\begin{proof}
\emph{(i)} By Proposition~\ref{prop:f5-n3crit} some $x_i \ge 2$. As $x_i > 0$,
the zero of $v_i$ lies at $\sigma(k)$ for some $k \ne i$. Reassigning by the
transposition of $i$ and $k$ changes $F$ by
$\big[v_i(\sigma(k)) + v_k(\sigma(i))\big] - \big[x_i + x_k\big] =
v_k(\sigma(i)) - x_i - x_k$, which is non-negative by optimality, so
$v_k(\sigma(i)) \ge x_i + x_k \ge 2$ and $\mathrm{sp}_k \ge 2$. Together with
$\mathrm{sp}_i \ge x_i \ge 2$ this gives $\Sigma(B) \ge 4$.

\emph{(ii)} By Proposition~\ref{prop:f5-n3crit}(b) there are distinct $i,j,k$
with $x_i = x_j = 1$, $v_i(\sigma(j)) = 0$ and $v_j(\sigma(k)) = 0$; here
$F(\sigma) = 2 + x_k$. Transposing $i$ and $j$ gives cost
$v_i(\sigma(j)) + v_j(\sigma(i)) + x_k = v_j(\sigma(i)) + x_k$, so optimality
forces $v_j(\sigma(i)) \ge 2$ and $\mathrm{sp}_j \ge 2$. Now take the
three-cycle $i \mapsto \sigma(j)$, $j \mapsto \sigma(k)$, $k \mapsto \sigma(i)$,
of cost
$v_i(\sigma(j)) + v_j(\sigma(k)) + v_k(\sigma(i)) = v_k(\sigma(i))$; optimality
forces $v_k(\sigma(i)) \ge 2 + x_k \ge 2$, so $\mathrm{sp}_k \ge 2$. With
$\mathrm{sp}_i \ge x_i = 1$ we get $\Sigma(B) \ge 1 + 2 + 2 = 5$.

The final claim is the contrapositive of (i) and (ii) together.
\end{proof}

Theorem~\ref{thm:f5-sigma3} uses \emph{only} optimality of the assignment. It
needs no minimality of the family, no balance hypothesis and no exchange step,
so it passes clear of all three closures of Remark~\ref{rem:f5-closures}.

\begin{corollary}
\label{cor:f5-contains-balanced}
Theorem~\ref{thm:f5-sigma3} contains Theorem~\ref{thm:balanced-class}.
\end{corollary}

\begin{proof}
A uniformly balanced family has $\mathrm{sp}_i \le 1$ for each of the three
agents, hence $\Sigma \le 3$.
\end{proof}

It is strictly wider: the spread profiles $(2,1,0)$ and $(3,0,0)$ also have
$\Sigma \le 3$ and are not uniformly balanced, so instances outside
\Cref{def:unifbal} are covered too --- including residual instances, on which no
uniformly balanced family exists at all.

\begin{corollary}
\label{cor:f5-reduction}
Conjecture~\ref{conj:f5} holds at $n = 3$ for every instance admitting a family
with $\Sigma \le 3$.
\end{corollary}

\begin{proof}
The $\Sigma$-minimal family of such an instance has $\Sigma \le 3$; apply
Theorem~\ref{thm:f5-sigma3}.
\end{proof}

So the whole of $n = 3$ rests on one purely combinatorial statement, with no
allocation, assignment or envy graph in it.

\begin{conjecture}[Total spread three --- \textbf{open}]
\label{conj:f5-sigma3}
Every three-agent negative dichotomous instance admits a partition of $\items$
into three bundles with $\sum_i \mathrm{sp}_i \le 3$.
\end{conjecture}

% --------------------------------------------------------------------------
\subsection{What stands in the way if Conjecture~\ref{conj:f5-sigma3} fails}
\label{sec:f5-pattern}

Theorem~\ref{thm:f5-sigma3} localises the obstruction completely at
$\Sigma = 4$, and the answer is a single rigid pattern.

\begin{proposition}[The unique obstruction at $\Sigma = 4$]
\label{prop:f5-pattern}
Let $n = 3$, let $\Sigma(B) = 4$ and let a minimum-cost $\sigma$ be bad. Then,
writing coordinates in the order $(\sigma(i),\sigma(j),\sigma(k))$ for a suitable
labelling of the agents,
\[
  v_i = v_k = (2,2,0), \qquad v_j = (0,0,0).
\]
\end{proposition}

\begin{proof}
By Theorem~\ref{thm:f5-sigma3}(ii) a two-path failure needs $\Sigma \ge 5$, so
the failure is an arc of weight at least $2$ and part~(i) applies: with $k$ the
holder of the zero of $v_i$, we have $\mathrm{sp}_i, \mathrm{sp}_k \ge 2$. Since
$\Sigma = 4$, this forces $\mathrm{sp}_i = \mathrm{sp}_k = 2$ and
$\mathrm{sp}_j = 0$, so $v_j \equiv 0$ and $x_j = 0$. Then $x_i \le
\mathrm{sp}_i = 2$ and $x_i \ge 2$ give $x_i = 2$, and
$v_k(\sigma(i)) \ge x_i + x_k = 2 + x_k$ with $v_k \le 2$ gives $x_k = 0$ and
$v_k(\sigma(i)) = 2$. Transposing $i$ and $j$ costs
$v_i(\sigma(j)) + v_j(\sigma(i)) + x_k = v_i(\sigma(j))$, so optimality gives
$v_i(\sigma(j)) \ge 2$, hence $= 2$. Finally the three-cycle
$i \mapsto \sigma(k)$, $j \mapsto \sigma(i)$, $k \mapsto \sigma(j)$ costs
$v_i(\sigma(k)) + v_j(\sigma(i)) + v_k(\sigma(j)) = v_k(\sigma(j))$, so
$v_k(\sigma(j)) \ge 2$, hence $= 2$. Collecting, $v_i = v_k = (2,2,0)$ and
$v_j = (0,0,0)$.
\end{proof}

Two agents therefore have \emph{identical} normalised vectors --- both see two
heavy bundles and one light one --- while the third is indifferent, and every
minimum-cost assignment strands one of the two on a heavy bundle. Ruling the
pattern out means re-partitioning the two heavy bundles so as to balance
\emph{both} of those agents at once, which is Conjecture~\ref{conj:f5-2balance}
below.

% --------------------------------------------------------------------------
\subsection{Two exchange lemmas, and one more conjecture}
\label{sec:f5-exchange}

\begin{lemma}[Intermediate value]
\label{lem:f5-ivt}
Let $\cost$ be dichotomous and $Y \subseteq \items$. For every integer
$0 \le k \le \cost(Y)$ there is $Z \subseteq Y$ with $\cost(Z) = k$.
\end{lemma}

\begin{proof}
Along a maximal chain $\emptyset = Z_0 \subset \dots \subset Z_r = Y$ consecutive
costs differ by $0$ or $1$, so the values $\cost(Z_0), \dots, \cost(Z_r)$ sweep
every integer between $0$ and $\cost(Y)$.
\end{proof}

Lemma~\ref{lem:f5-ivt} is what replaces the exchange step forbidden by the third
closure of Remark~\ref{rem:f5-closures}: single items need not move a cost, but
\emph{sets} realising any intermediate value always exist.

\begin{lemma}[Two-way balance, one agent]
\label{lem:f5-balance1}
For dichotomous $\cost$ and any $U \subseteq \items$ there is $Z \subseteq U$
with $\abs{\cost(Z) - \cost(U \setminus Z)} \le 1$.
\end{lemma}

\begin{proof}
On a maximal chain $\emptyset = Z_0 \subset \dots \subset Z_r = U$ put
$f(t) := \cost(Z_t) - \cost(U \setminus Z_t)$. Then $f(0) = -\cost(U) \le 0$ and
$f(r) = \cost(U) \ge 0$. Each step raises the first term by $0$ or $1$ and
lowers the second by $0$ or $1$, so $f$ is non-decreasing with steps in
$\set{0,1,2}$. A step of $2$ cannot carry $f$ from at most $-2$ to at least $2$,
so some $f(t)$ lies in $\set{-1,0,1}$.
\end{proof}

\begin{conjecture}[Two-way balance, two agents --- \textbf{open}]
\label{conj:f5-2balance}
For any two dichotomous costs and any $U \subseteq \items$ there is a partition
$U = Z \sqcup (U \setminus Z)$ with $\mathrm{sp}_1 \le 1$ and
$\mathrm{sp}_2 \le 1$ on the two parts.
\end{conjecture}

Note that Proposition~\ref{prop:no-balance} obstructs the \emph{three}-agent
analogue of Conjecture~\ref{conj:f5-2balance}, not this one, so no barrier is
known.

\begin{remark}[The two-agent statement is a synchronisation problem]
\label{rem:f5-walk}
For $Z \subseteq U$ put
\[
  g(Z) \;:=\; \Big(\cost_1(Z) - \cost_1(U \setminus Z),\;
                   \cost_2(Z) - \cost_2(U \setminus Z)\Big),
\]
so that Conjecture~\ref{conj:f5-2balance} asks for $g(Z) \in \set{-1,0,1}^2$;
note $g(U \setminus Z) = -g(Z)$, so the image is symmetric about the origin.
Along any maximal chain both coordinates of $g$ are \emph{non-decreasing} with
steps in $\set{0,1,2}$, running from $(-\cost_1(U), -\cost_2(U))$ to
$(\cost_1(U), \cost_2(U))$. Lemma~\ref{lem:f5-balance1} is the one-coordinate
case of this. The two-coordinate case does not follow, because the coordinates
cross the band $\set{-1,0,1}$ at different times: at the first $t$ with
$g_2(t) \ge -1$ one gets $g_2(t) \in \set{-1,0}$, but $g_1(t)$ may already be
far above $1$. The content of Conjecture~\ref{conj:f5-2balance} is exactly that
the \emph{order} of the chain --- the only freedom available --- can be chosen to
synchronise the two crossings. For additive measures this is the discrete
necklace-splitting theorem with two cuts; dichotomous costs are not additive, so
that result does not transfer.
\end{remark}

% --------------------------------------------------------------------------
\subsection{Status}
\label{sec:f5-status}

At $n = 3$ everything above is proved except two statements, and either one
closes the case:

\begin{center}\small
\begin{tabular}{@{}ll@{}}
\toprule
statement & status \\
\midrule
Observation~\ref{obs:f5-arcs}, Proposition~\ref{prop:f5-n3crit} & proved \\
Theorem~\ref{thm:f5-sigma3} ($\Sigma \le 3$ suffices) & proved \\
Corollary~\ref{cor:f5-contains-balanced} (contains Thm~\ref{thm:balanced-class})
  & proved \\
Corollary~\ref{cor:f5-reduction} (the reduction) & proved \\
Proposition~\ref{prop:f5-pattern} (the unique $\Sigma = 4$ pattern) & proved \\
Lemma~\ref{lem:f5-ivt}, Lemma~\ref{lem:f5-balance1} & proved \\
\textbf{Conjecture~\ref{conj:f5-sigma3}} ($\min \Sigma \le 3$) & \textbf{open} \\
\textbf{Conjecture~\ref{conj:f5-2balance}} (two-agent balance) & \textbf{open} \\
\bottomrule
\end{tabular}
\end{center}

\begin{remark}[Why Conjecture~\ref{conj:f5-sigma3} is plausible, and how it would
fail]
\label{rem:f5-sigma3-odds}
On a residual instance no family has all three spreads at most $1$, so
$\Sigma \le 3$ must be attained as $(2,1,0)$, $(2,0,0)$ or $(3,0,0)$ --- each of
which needs some agent to have spread exactly $0$, that is, to value all three
bundles equally. For a binary additive agent that means the three bundles hold
equally many elements of $D_i$, which is impossible unless $3$ divides
$\abs{D_i}$. So a counterexample needs three binary additive agents with no
$\abs{D_i}$ divisible by $3$ and no uniformly balanced family.

The two requirements pull against each other, which is the reason to expect the
conjecture to hold. Writing $\abs{D_i} = 3q_i + r_i$, uniform balance asks the
counts of $D_i$ to be $(q_i{+}1,q_i,q_i)$ or $(q_i{+}1,q_i{+}1,q_i)$ up to
order when $r_i \ne 0$ --- in each case a free choice of which bundle is heavy.
Only $3 \mid \abs{D_i}$ makes the count vector rigid, and it is exactly that
rigidity which Proposition~\ref{prop:no-balance} exploits: its two size-three
sets are forced rainbow and its size-two set then contradicts them. Barring
sizes divisible by $3$ removes the rigidity that makes uniform balance fail,
while being precisely what is needed to block spread $0$.

This is a heuristic and it concerns binary additive agents only; a general
dichotomous agent can reach spread $0$ with no divisibility condition, as
$\cost(S) = \min(\abs{S},1)$ does on any partition into non-empty bundles.
\end{remark}

\paragraph{Verdict.} Conjecture~\ref{conj:f5} is open, unrefuted, constructive
and always applicable, which no earlier target in this report has been at once.
At $n = 3$ it is reduced to Conjecture~\ref{conj:f5-sigma3}, a statement about
partitions alone. Nothing here proves any case with $n \ge 3$, and the reduction
should not be represented as one.
